% Generated draft: requires mathematical and source review.
\documentclass[11pt]{article}
\usepackage[margin=25mm]{geometry}
\usepackage{fontspec}
\setmainfont{Latin Modern Roman}
\usepackage{amsmath,amssymb,mathtools,unicode-math}
\setmathfont{Latin Modern Math}
\usepackage{enumitem,needspace,xurl,hyperref}
\setlength{\parindent}{0pt}
\setlength{\parskip}{5pt}
\setlength{\emergencystretch}{3em}
\newcommand{\N}{\mathbb{N}}
\newcommand{\Z}{\mathbb{Z}}
\newcommand{\Q}{\mathbb{Q}}
\newcommand{\R}{\mathbb{R}}
\newcommand{\C}{\mathbb{C}}
\newcommand{\F}{\mathbb{F}}
\newcommand{\A}{\mathbb{A}}
\newcommand{\PP}{\mathbb{P}}
\newcommand{\E}{\mathbb{E}}
\newcommand{\eps}{\varepsilon}
\newcommand{\dd}{\,\mathrm{d}}
\begin{document}
\section*{UnsolvedMath ID: 30003119}
\textbf{Source problem code:} catalogue entry
\subsection{Definitions and mathematical statement}
Let \(A\) be a complex algebraic variety, let \(f:A\to A\) be an algebraic endomorphism satisfying the source condition \(\deg(f)>1\), let \(V\subseteq A\) be a subvariety, and let \(P\in A(\C)\). Define \(f^0=\operatorname{id}_A\) and \(f^{n+1}=f\circ f^n\). Using \(\Z_{\geq 0}\) for the iteration indices, define the hitting-time set
\[
S_f(P,V)=\{n\in\Z_{\geq 0}:f^n(P)\in V(\C)\}.
\]
The supplied question asks whether, for every such \(A,f,V,P\), the set \(S_f(P,V)\) is a finite union of one-sided arithmetic progressions. If progressions of step zero are allowed, this conclusion means that there exist \(r\in\Z_{\geq 0}\) and integers \(a_i,b_i\geq 0\), for \(1\leq i\leq r\), such that
\[
S_f(P,V)=\bigcup_{i=1}^{r}\{a_i+b_i k:k\in\Z_{\geq 0}\}.
\]
Here \(r=0\) represents the empty set, and \(b_i=0\) represents a singleton. Equivalently, this convention permits a finite exceptional set together with finitely many progressions of positive step.

\textbf{Unresolved qualifications.} The source does not specify the meaning of \(\deg(f)\) for its class of varieties and endomorphisms. For an integral variety and a dominant generically finite endomorphism, one possible meaning is the function-field degree \([\C(A):f^*\C(A)]\); the source does not explicitly supply these hypotheses or select this degree convention. Thus the degree condition is retained pending clarification, rather than used to impose additional hypotheses. The source also does not specify whether varieties must be irreducible, whether subvarieties must be closed, or whether additional geometric assumptions are intended.

The displayed formulation explicitly uses nonnegative iteration indices and permits singleton progressions. The source writes \(\N\) without specifying whether zero is included, and does not define its arithmetic-progression convention. If only positive steps are allowed, the displayed conclusion has an additional allowance for isolated times that requires confirmation. With singleton progressions permitted, changing whether the initial time zero is included does not change the finite-union property.
\subsection{Short English statement}
The question concerns the times when an iterated algebraic endomorphism of a complex variety sends a given point into a given subvariety. Must these times form finitely many one-sided arithmetic progressions, allowing isolated times under the convention stated above? The degree hypothesis and several geometric conventions require clarification from the source.
\subsection{Sources}
\begin{itemize}
\item Local UnsolvedMath registry, supplied complete record: permanent ID \(30003119\), exact source problem code \texttt{catalogue entry}, title ``The Dynamical Mordell--Lang Conjecture for Endomorphisms,'' in the Oberwolfach Reports Open Problems collection. The supplied record gives proposed year 2016 and names no individual proposer.
\item Supplied report citation: \emph{Diophantische Approximationen} (2017), Oberwolfach Reports, pp.\ 1099--1187, DOI \texttt{10.4171/owr/2016/21}. \url{https://doi.org/10.4171/owr/2016/21}. The supplied background describes these reports as documenting mathematical workshops and their open-problem discussions.
\item Supplied literature source, ``OWR report'': \url{https://ems.press/content/serial-article-files/46625?nt=1}.
\item Supplied literature source, ``Recent dynamical Mordell--Lang work'': \url{https://academic.oup.com/imrn/article-abstract/2025/9/rnaf115/8126143}.
\item Literature assessment dated exactly \texttt{2026-08-21}: the supplied curation labels the problem partially solved and states that special cases are proved, that the general complex characteristic-zero statement in this record remains open, and that recent positive-characteristic advances do not settle that general case. This is supplied dated curation derived from the OWR material, not an independently verified assessment; its time-sensitive status and cited claims require independent verification.
\item The registry labels the statement \texttt{corrected\_verified}, explaining that mathematical content was preserved while notation, typography, or omitted definitions were repaired. This is supplied registry metadata, not independent mathematical verification of this draft.
\end{itemize}
\end{document}
